
\section{Identificación}

\begin{tabularx}{\textwidth}{@{}lX@{}}
\toprule
Campo & Valor \\
\midrule
Nombre del CADi & IA agéntica para la creación de contenido académico \\
Escuela / unidad académica & Escuela de Ingeniería y Ciencias (EIC) \\
Clave sugerida & \clave\ \textit{(placeholder; asignar clave oficial con CEDDIE / EIC)} \\
Modalidad & Virtual (Microsoft Teams) / híbrida / presencial según edición \\
Duración & \textbf{40 horas totales}: 24\,h en aula (8 sesiones $\times$ 3\,h) + 16\,h fuera de aula \\
Cupo sugerido & 20--30 participantes \\
Idioma & Español (México) \\
Nivel / audiencia & Profesorado de profesional y posgrado de la EIC \\
Instructor & Juliho Castillo \textit{(completar CV institucional, correo y perfil)} \\
Unidad / área de apoyo & CEDDIE / Innovación Educativa · oferta de CADi de la EIC \\
\bottomrule
\end{tabularx}

\begin{quote}
Formato orientativo tipo CEDDIE. Campos y cupos son \textbf{sugerencias}; \ceddie\ antes de la oferta oficial.
\end{quote}

\section{Justificación}

En la \textbf{Escuela de Ingeniería y Ciencias}, el profesorado produce de forma continua materiales típicos de STEM: \textbf{textbooks y ebooks técnicos}, \textbf{guías de laboratorio}, \textbf{presentaciones de clase y de conferencia}, y \textbf{artículos de investigación o de divulgación técnica}. Estos artefactos exigen coherencia en \textbf{notación, unidades, terminología y objetivos de aprendizaje}.

Las herramientas de IA generativa ya se usan de forma informal (chat único, prompts ad hoc), pero ese uso suele ser \textbf{fragmentario}: se pierde contexto entre sesiones, no hay roles claros, y es difícil auditar qué hizo la máquina y qué decidió la persona.

La \textbf{IA agéntica} introduce un cambio de paradigma: se diseñan \textbf{agentes} con roles (planificador, redactor, revisor, verificador de consistencia), \textbf{herramientas} y \textbf{orquestación} con supervisión humana (\emph{human-in-the-loop}). Esto permite:

\begin{itemize}
\item Acelerar el paso de idea $\rightarrow$ outline $\rightarrow$ borrador $\rightarrow$ revisión en obras técnicas largas.
\item Mantener un \textbf{registro auditable} (agent log) útil para integridad académica.
\item Estandarizar plantillas reutilizables entre cursos EIC.
\item Reducir inconsistencias de notación, unidades, glosario y alineación didáctica.
\end{itemize}

El formato de \textbf{40 horas} (24 en aula + 16 independientes) permite profundizar en orquestación, dividir cada tipología de contenido en dos sesiones y sostener avance real del \textbf{microartefacto}.

Cuando la IA interviene en \textbf{ecuaciones, citas o afirmaciones basadas en datos}, la verificación humana es crítica.

\begin{quote}
\textbf{Verificar con CEDDIE / política institucional vigente} los lineamientos de integridad académica, uso de IA generativa, protección de datos y propiedad intelectual aplicables.
\end{quote}

\section{Perfil del participante}

Profesoras y profesores de profesional o posgrado de la EIC que:

\begin{itemize}
\item Diseñan o actualizan materiales STEM.
\item Ya exploraron IA generativa a nivel conversacional y desean pasar a flujos estructurados.
\item Requieren criterios claros de calidad, autoría y transparencia.
\item Pueden dedicar trabajo independiente entre sesiones.
\end{itemize}

\textbf{No es un curso de programación.} El énfasis es el diseño de flujos agénticos y la calidad del contenido académico.

\section{Prerrequisitos}

\begin{enumerate}
\item Experiencia docente activa en profesional o posgrado (preferentemente EIC).
\item Manejo básico de un asistente de IA generativa (chat).
\item Cuenta institucional y acceso a la plataforma de videoconferencia de la edición.
\item Recomendado: un tema o material propio EIC sobre el cual trabajar.
\item Disponibilidad para $\sim$2\,h de trabajo independiente por sesión (16\,h totales fuera de aula).
\end{enumerate}

\section{Resultados de aprendizaje}

Al finalizar el CADi, la persona participante:

\begin{tabularx}{\textwidth}{@{}cX@{}}
\toprule
\# & Resultado de aprendizaje \\
\midrule
RA1 & Explica la diferencia entre chat conversacional y sistemas agénticos (roles, herramientas, orquestación, HITL) y argumenta su relevancia en materiales EIC. \\
RA2 & Elabora un brief de contenido, un prompt de sistema y un esquema de orquestación multi-rol reutilizable. \\
RA3 & Produce, con supervisión humana, la arquitectura y outline de un libro/ebook técnico alineado a resultados de aprendizaje. \\
RA4 & Genera secciones, ejercicios/figuras y un pase de consistencia (notación, unidades, terminología). \\
RA5 & Diseña storyboard y narrativa de una presentación didáctica alineada a objetivos. \\
RA6 & Construye un deck mínimo viable con notas de orador y criterios de accesibilidad. \\
RA7 & Estructura y redacta (con supervisión) un micro-artículo o sección académica (p.\ ej.\ IMRyD), documentando citas y verificación. \\
RA8 & Aplica agentes de revisión, completa agent log + checklist y presenta un microartefacto EIC como evidencia de cierre. \\
\bottomrule
\end{tabularx}

\section{Contenido temático}

\subsection{En aula (24\,h --- 8 sesiones $\times$ 3\,h)}

\begin{enumerate}
\item Fundamentos agénticos e integridad académica (EIC).
\item Diseño de briefs, prompts de sistema y orquestación.
\item Libros/ebooks I --- arquitectura de obra y outline.
\item Libros/ebooks II --- secciones, ejercicios, consistencia.
\item Presentaciones I --- objetivos, narrativa, storyboard.
\item Presentaciones II --- deck, notas de orador, accesibilidad.
\item Artículos I --- estructura, borrador asistido, citas.
\item Artículos II --- revisión agéntica, cierre, showcase.
\end{enumerate}

\subsection{Fuera de aula (16\,h)}

\begin{tabularx}{\textwidth}{@{}Xccc@{}}
\toprule
Bloque & Horas & Apoya & Descripción \\
\midrule
Lecturas y demos guiadas & 4\,h & 1--2 & Lecturas curadas, demos del stack \\
Práctica con stack y plantillas & 4\,h & 2--6 & Ensayo de plantillas \\
Avance del microartefacto & 6\,h & 3--8 & Desarrollo progresivo del entregable EIC \\
Agent log + checklist + reflexión & 2\,h & 7--8 & Documentación auditable \\
\textbf{Total} & \textbf{16\,h} & --- & --- \\
\bottomrule
\end{tabularx}

\section{Metodología}

\begin{itemize}
\item Taller práctico con demostraciones en vivo y trabajo individual/parejas sobre material propio EIC.
\item Aprendizaje basado en proyecto: flujo agéntico y microartefacto final con agent log.
\item Trabajo independiente explícito (tabla de 16\,h).
\item Plantillas (brief, prompt de sistema, checklist) como artefactos transferibles.
\item Reflexión ética integrada: transparencia, verificación de fuentes, límites de la automatización.
\end{itemize}

\section{Evaluación y evidencia de acreditación}

\begin{tabularx}{\textwidth}{@{}XccX@{}}
\toprule
Componente & Peso & Evidencia \\
\midrule
Participación y actividades en aula & 25\,\% & Entregables E1--E7 \\
Trabajo independiente (16\,h) & 15\,\% & Lecturas/práctica y avance documentado \\
Uso de plantillas & 15\,\% & Archivos completados y reutilizados \\
Microartefacto final & 30\,\% & Capítulo STEM / deck técnico / IMRyD / guía de lab \\
Agent log + reflexión de integridad & 15\,\% & Registro de pasos y verificación \\
\bottomrule
\end{tabularx}

\textbf{Criterio de acreditación sugerido:} cumplir el umbral definido por CEDDIE para CADi (\ceddie) y entregar el microartefacto con agent log y checklist.

\section{Recursos y herramientas}

Enfoque \textbf{agnóstico a proveedor}. Stack típico: plataforma de videoconferencia, asistente(s) de IA con instrucciones de sistema, editor colaborativo; opcional: diapositivas, gestores de citas, editores de ecuaciones.

\begin{quote}
\ceddie\ qué herramientas están homologadas, restricciones de datos sensibles y uso de cuentas institucionales.
\end{quote}

\section{Perfil del instructor}

\begin{tabularx}{\textwidth}{@{}lX@{}}
\toprule
Campo & Contenido \\
\midrule
Nombre & Juliho Castillo \\
Rol & Instructor / facilitador del CADi \\
Escuela / unidad & Escuela de Ingeniería y Ciencias (EIC) \\
Experiencia relevante & \textit{(completar)} \\
Correo institucional & \textit{(completar)} \\
Campus / escuela & \textit{(completar; EIC)} \\
Bio corta & \textit{(completar 80--120 palabras)} \\
Co-facilitadores & \textit{(opcional; CEDDIE / EIC)} \\
\bottomrule
\end{tabularx}

\section{Observaciones para la convocatoria}

\begin{itemize}
\item Enfatizar que el CADi es práctico y transferible, orientado al profesorado de la EIC.
\item Comunicar claramente: \textbf{40\,h (24 aula + 16 independientes)}.
\item Pedir tema propio EIC (capítulo STEM, guía de lab, deck técnico o sección de artículo).
\item Indicar políticas de integridad y de uso de IA (\ceddie).
\item Mencionar colaboración con CEDDIE / oferta de CADi de la Escuela.
\item Actualizar clave, cupo, fechas y stack tecnológico por edición.
\end{itemize}
